% !TeX root = ../main.tex
% =========================================================================
% PREÁMBULO - Formato Institucional UPS Sede Guayaquil (APA 7ma Edición)
% =========================================================================

% --- Idioma y codificación ---
\usepackage[utf8]{inputenc}
\usepackage[spanish,mexico,es-noshorthands]{babel}

% Comando para obtener el nombre del mes actual en español
\newcommand{\monthname}{%
  \ifcase\month\or enero\or febrero\or marzo\or abril\or mayo\or junio\or julio\or agosto\or septiembre\or octubre\or noviembre\or diciembre\fi}

% --- Tipografía APA 7 (Times New Roman 12pt en todo el documento) ---
\usepackage{newtxtext,newtxmath} % Times New Roman para texto y matemáticas
\usepackage[T1]{fontenc}
\usepackage{anyfontsize}
\usepackage{amsmath}
\usepackage{indentfirst} % Sangría APA 7 en TODOS los párrafos, incluido el primero tras una sección
\usepackage{enumitem}    % Sangría APA 7 en listas con viñetas y numeradas
\setlist[itemize]{leftmargin=1.27cm, labelsep=0.3cm, topsep=4pt, itemsep=2pt}
\setlist[enumerate]{leftmargin=1.27cm, labelsep=0.3cm, topsep=4pt, itemsep=2pt}
\usepackage{float}
\usepackage{pdflscape}
\usepackage[
    a4paper,
    left=2.5cm,
    right=2.5cm,
    top=2.5cm,
    bottom=2.5cm
]{geometry}

% --- Gráficos y Colores ---
\usepackage[table,xcdraw]{xcolor} % Carga xcolor una sola vez (incluye colortbl)
\usepackage{graphicx}
\usepackage{tikz}
\usetikzlibrary{calc,arrows,arrows.meta,positioning,shapes.geometric,fit}
\usepackage{caption}

% --- Tablas y Organización ---
\usepackage{booktabs}
\usepackage{array}
\usepackage{longtable}
\usepackage{multicol}

% --- Parche de compatibilidad (solo actúa si hace falta) ---
% En instalaciones locales de MiKTeX con array.sty antiguo y colortbl.sty
% nuevo, las columnas p{} fallan con "Undefined control sequence \insert@pcolumn".
% En Overleaf o con MiKTeX actualizado este bloque no hace nada.
\makeatletter
\@ifundefined{insert@pcolumn}{\let\insert@pcolumn\insert@column}{}
\makeatother

% --- Formato APA 7 de Párrafos ---
\setlength{\parindent}{1.27cm} % Sangría de primera línea APA 7 (0.5 in)
\setlength{\parskip}{0pt}      % Sin espacio adicional entre párrafos

% Nota de pie para tablas y figuras en formato APA 7
\newcommand{\notatabla}[1]{\par\vspace{4pt}\raggedright\footnotesize\textit{Nota.} #1}

% --- Encabezados y pies (APA 7: número de página en la esquina superior derecha) ---
\usepackage{fancyhdr}
\pagestyle{fancy}
\fancyhf{}
\setlength{\headheight}{15pt}
\addtolength{\topmargin}{-3pt}
\rhead{\thepage}
\renewcommand{\headrulewidth}{0pt}

% --- Configuración de Secciones (Todas a 12pt) ---
\usepackage{titlesec}
\titleformat{\section}
  {\normalfont\fontsize{12}{14}\bfseries}
  {\thesection.}
  {1em}
  {}
\titleformat{\subsection}
  {\normalfont\fontsize{12}{14}\bfseries}
  {\thesubsection.}
  {1em}
  {}
\titleformat{\subsubsection}
  {\normalfont\fontsize{12}{14}\bfseries\itshape}
  {\thesubsubsection.}
  {1em}
  {}

% --- Índices a 12pt ---
\providecommand{\UseTaggingSocket}[1]{}
\providecommand{\SuspendTagging}[1]{}
\providecommand{\ResumeTagging}[1]{}
\usepackage{tocloft}
\renewcommand{\cfttoctitlefont}{\normalfont\fontsize{12}{14}\bfseries}
\renewcommand{\cftsecfont}{\normalfont\fontsize{12}{14}\selectfont}
\renewcommand{\cftsecpagefont}{\normalfont\fontsize{12}{14}\selectfont}
\renewcommand{\cftsubsecfont}{\normalfont\fontsize{12}{14}\selectfont}
\renewcommand{\cftsubsecpagefont}{\normalfont\fontsize{12}{14}\selectfont}
\renewcommand{\cftsubsubsecfont}{\normalfont\fontsize{12}{14}\selectfont}
\renewcommand{\cftsubsubsecpagefont}{\normalfont\fontsize{12}{14}\selectfont}
\renewcommand{\cftsecdotsep}{\cftdotsep}
\renewcommand{\cftloftitlefont}{\normalfont\fontsize{12}{14}\bfseries}
\renewcommand{\cftfigfont}{\normalfont\fontsize{12}{14}\selectfont}
\renewcommand{\cftfigpagefont}{\normalfont\fontsize{12}{14}\selectfont}
\renewcommand{\cftlottitlefont}{\normalfont\fontsize{12}{14}\bfseries}
\renewcommand{\cfttabfont}{\normalfont\fontsize{12}{14}\selectfont}
\renewcommand{\cfttabpagefont}{\normalfont\fontsize{12}{14}\selectfont}

% --- Títulos de figuras y tablas en formato APA 7 ---
\captionsetup[table]{
    position=above,
    labelfont=bf,
    textfont=it,
    labelsep=newline,
    justification=justified,
    singlelinecheck=false,
    skip=6pt
}
\captionsetup[figure]{
    position=above,
    labelfont=bf,
    textfont=it,
    labelsep=newline,
    justification=justified,
    singlelinecheck=false,
    skip=6pt
}

% --- Interlineado 1.5 ---
\usepackage{setspace}
\onehalfspacing

% --- Bibliografía (APA 7 Estándar) ---
\usepackage{csquotes}
\usepackage[style=apa, backend=biber]{biblatex}
\setlength{\bibhang}{1.27cm} % Sangría francesa reglamentaria APA 7 (0.5 in)
\addbibresource{referencias/references.bib}

% --- Enlaces ---
\usepackage{hyperref}
\hypersetup{
    colorlinks=true,
    linkcolor=black,
    citecolor=blue,
    urlcolor=blue
}

% --- Citas en azul con enlace (parche de la plantilla oficial) ---
\DeclareFieldFormat{citehyperref}{\bibhyperref{\textcolor{blue}{#1}}}
\DeclareNameFormat{labelname:hyperref}{%
  \bibhyperref{\textcolor{blue}{%
    \ifcase\value{uniquename}%
      \usebibmacro{name:family}{\namepartfamily}{\namepartgiven}{\namepartprefix}{\namepartsuffix}%
    \or
      \usebibmacro{name:given-family}{\namepartfamily}{\namepartgiveni}{\namepartprefix}{\namepartsuffix}%
    \or
      \usebibmacro{name:given-family}{\namepartfamily}{\namepartgiven}{\namepartprefix}{\namepartsuffix}%
    \fi}}%
  \usebibmacro{name:andothers}}

\makeatletter
\DeclareCiteCommand{\parencite}[\mkbibparens]
  {\usebibmacro{prenote}}
  {\printtext[citehyperref]{%
     \printnames{labelname}%
     \setunit{\printdelim{nameyeardelim}}%
     \printlabeldateextra}}
  {\multicitedelim}
  {\usebibmacro{postnote}}

\DeclareCiteCommand{\textcite}
  {\boolfalse{cbx:parens}}
  {\printtext[citehyperref]{\printnames{labelname}}%
   \setunit{\addspace}%
   \printtext[citehyperref]{%
     \ifnumequal{\value{citecount}}{1}{\bibopenparen}{}%
     \printlabeldateextra}}
  {\ifbool{cbx:parens}{\bibcloseparen\global\boolfalse{cbx:parens}}{}\multicitedelim}
  {\ifnumequal{\value{citecount}}{\value{citetotal}}{\bibcloseparen}{}\usebibmacro{postnote}}
\makeatother

% --- URLs dentro de márgenes ---
\usepackage{url}
\expandafter\def\expandafter\UrlBreaks\expandafter{\UrlBreaks%
  \do\/\do\-\do\.\do\0\do\1\do\2\do\3\do\4\do\5\do\6\do\7\do\8\do\9\do\a\do\b\do\c\do\d\do\e\do\f\do\g\do\h\do\i\do\j\do\k\do\l\do\m\do\n\do\o\do\p\do\q\do\r\do\s\do\t\do\u\do\v\do\w\do\x\do\y\do\z\do\A\do\B\do\C\do\D\do\E\do\F\do\G\do\H\do\I\do\J\do\K\do\L\do\M\do\N\do\O\do\P\do\Q\do\R\do\S\do\T\do\U\do\V\do\W\do\X\do\Y\do\Z}
\urlstyle{same}

% --- Datos del proyecto (se editan en un solo lugar) ---
\newcommand{\tituloProyecto}{Desarrollo de un agente inteligente para la atención comercial y cotización de materiales eléctricos en Guayaquil}
\newcommand{\autorUnoProyecto}{Victor Jesus Totoy Iza}
\newcommand{\autorDosProyecto}{Rodolfo Damián Seminario Murillo}
\newcommand{\autoresProyecto}{Victor Jesus Totoy Iza y Rodolfo Damián Seminario Murillo}
\newcommand{\tutorProyecto}{Ing. Joe Llerena Izquierdo, M.Sc., Ph.D.}
